\documentclass[11pt]{article}
\usepackage{amsmath,amssymb,amsthm}
\usepackage[colorlinks=true]{hyperref}
\usepackage[margin=1in]{geometry}
\usepackage{booktabs}
\usepackage{array}

\newtheorem{theorem}{Theorem}
\newtheorem{lemma}[theorem]{Lemma}
\newtheorem{proposition}[theorem]{Proposition}
\newtheorem{corollary}[theorem]{Corollary}
\newtheorem{conjecture}[theorem]{Conjecture}
\newtheorem{observation}[theorem]{Observation}
\newtheorem{question}[theorem]{Open Question}
\theoremstyle{definition}
\newtheorem{definition}[theorem]{Definition}
\theoremstyle{remark}
\newtheorem{remark}[theorem]{Remark}

\title{Sequence Memory Capacity in Small Recurrent Networks\\--- Research Session 2026-03-26}
\author{Research Notebook}
\date{March 26, 2026}

\begin{document}
\maketitle

\section{Objective}

Given a recurrent network of $N \leq 100$ neurons that observes a sequence of patterns \emph{once}, what is the maximum sequence length $L(N)$ that can be perfectly replayed? We seek:
\begin{enumerate}
    \item A practical framework: working implementations with empirical capacity measurements across architectures and coding distributions.
    \item Theoretical results: capacity bounds as a function of $N$, the feature dimension $M$, and the pattern distribution.
\end{enumerate}

The investigation was motivated by the observation that biological neurons exhibit Laplace-distributed firing frequencies, which is entropy-maximizing under metabolic (L1) constraints. We explore whether this has implications for optimal sequence storage.

\section{Approach}

We studied three architectures for one-shot sequence memorization, each with a different capacity regime:

\begin{itemize}
    \item \textbf{Architecture A (Linear Pseudoinverse):} $\mathbf{x}_{k+1} = W\mathbf{x}_k$ with $W = X_2 X_1^+$ (Moore--Penrose pseudoinverse). The simplest possible model.
    \item \textbf{Architecture B (Echo State / Random Features):} $\mathbf{x}_{k+1} = W_2 \sigma(W_1 \mathbf{x}_k)$ where $W_1 \in \mathbb{R}^{M \times N}$ is a fixed random projection, $\sigma$ is a pointwise nonlinearity (tanh), and $W_2 = X_2 [\sigma(W_1 X_1)]^+$.
    \item \textbf{Architecture C (Modern Hopfield / Softmax Retrieval):} $\mathbf{x}_{k+1} = X_2 \operatorname{softmax}(\beta X_1^\top \mathbf{x}_k)$, storing key-value pairs $(X_1, X_2)$ explicitly.
\end{itemize}

Four pattern distributions were compared: Gaussian, Laplace, sparse (10\% nonzero), and orthogonal, all with matched variance.

The session proceeded as: literature survey (Wide Mode, 13:49--13:55), then theory development and simulation (Deep Mode, 13:55--15:50).

\section{Background \& Prior Work}

\subsection{Classical Sequence Memory}

The capacity of associative memories for \emph{static} patterns is well-studied. Classical Hopfield networks store $\sim 0.14N$ patterns (Hopfield 1982, Amit et al.\ 1987). For \emph{sequences}, Sompolinsky and Kanter (1986) showed capacity $\sim 0.27N$ with binary neurons and Hebbian learning.

\subsection{One-Shot Sequence Storage}

A recent hippocampus-inspired model achieves one-shot sequence storage with capacity $\sim N$ patterns using Hebbian descent learning, pattern separation via a dentate gyrus module, and intrinsic CA3 dynamics~\cite{plos2024oneshot}. The key innovation is that pattern separation (reducing inter-pattern correlations) enables reliable hetero-association.

\subsection{Modern Hopfield Networks}

Ramsauer et al.\ (2020) showed that Hopfield networks with continuous states and a softmax (exponential) energy function achieve \emph{exponential} storage capacity: up to $2^{cN}$ patterns can be stored and retrieved in a single update step~\cite{ramsauer2020hopfield}. The update rule is equivalent to transformer key-value attention. A 2024 NeurIPS paper further tightened the capacity bounds, showing provably optimal results~\cite{neurips2024hopfield}.

\subsection{Matrix RNN Capacity}

For recurrent networks with matrix-valued states ($N \times N$ state matrices instead of $N$-dimensional vectors), the memory capacity scales as $N^2$, a quadratic improvement over vector networks~\cite{matrix_rnn2021}.

\subsection{Laplace Distribution in Neural Coding}

In computational neuroscience, sparse coding models use $L_1$ penalties (equivalent to Laplace priors) to produce efficient neural representations~\cite{sparse_coding}. The Laplace distribution is the maximum-entropy distribution under an $L_1$ constraint, making it information-theoretically optimal for neurons operating under metabolic (mean absolute activity) budgets.

\section{Findings}

\subsection{Theoretical Results}

\begin{definition}[One-shot sequence memory]
A \emph{one-shot sequence memory} is a system with $N$-dimensional state that, given a sequence $S = (\mathbf{x}_1, \ldots, \mathbf{x}_L)$ with $\mathbf{x}_k \in \mathbb{R}^N$, computes parameters $\theta(S)$ in a single pass, and then produces a map $F_\theta: \mathbb{R}^N \to \mathbb{R}^N$ such that $F_\theta(\mathbf{x}_k) = \mathbf{x}_{k+1}$ for all $k = 1, \ldots, L-1$.
\end{definition}

\begin{theorem}[Linear Capacity Bound]\label{thm:linear}
Let $F_\theta(\mathbf{x}) = W\mathbf{x}$ with $W \in \mathbb{R}^{N \times N}$ computed as $W = X_2 X_1^+$, where $X_1 = [\mathbf{x}_1, \ldots, \mathbf{x}_{L-1}]$ and $X_2 = [\mathbf{x}_2, \ldots, \mathbf{x}_L]$. Then:
\begin{enumerate}
    \item[(i)] Perfect recall ($W\mathbf{x}_k = \mathbf{x}_{k+1}$ for all $k$) holds if and only if $X_1$ has full column rank, which requires $L - 1 \leq N$.
    \item[(ii)] For patterns drawn i.i.d.\ from any absolutely continuous distribution on $\mathbb{R}^N$, the maximum achievable sequence length is $L^* = N + 1$ almost surely.
\end{enumerate}
\end{theorem}

\begin{proof}
The pseudoinverse satisfies $X_1^+ X_1 = I_{L-1}$ if and only if $X_1$ has full column rank, i.e., $\operatorname{rank}(X_1) = L - 1$. Since $X_1 \in \mathbb{R}^{N \times (L-1)}$, this requires $L - 1 \leq N$. When this holds, $WX_1 = X_2 X_1^+ X_1 = X_2$, which is exactly the perfect-recall condition.

For the achievability, any $N$ vectors drawn from an absolutely continuous distribution in $\mathbb{R}^N$ are linearly independent with probability 1. Setting $L - 1 = N$ gives $L^* = N + 1$.
\end{proof}

\begin{theorem}[Feature Expansion Capacity]\label{thm:feature}
Let $\varphi: \mathbb{R}^N \to \mathbb{R}^M$ be a feature map (e.g., $\varphi(\mathbf{x}) = \sigma(W_1 \mathbf{x})$ for fixed $W_1$). Define $F_\theta(\mathbf{x}) = W_2 \varphi(\mathbf{x})$ with $W_2 = X_2 \Phi^+$ where $\Phi = [\varphi(\mathbf{x}_1), \ldots, \varphi(\mathbf{x}_{L-1})]$. Then:
\begin{enumerate}
    \item[(i)] Perfect recall holds if and only if $\Phi$ has full column rank, requiring $L - 1 \leq M$.
    \item[(ii)] The maximum achievable sequence length is $L^* = M + 1$, provided $\varphi$ does not collapse distinct patterns (i.e., the features are in general position).
    \item[(iii)] The total parameter count is $P = NM + MN = 2NM$ (for $W_1$ and $W_2$), giving capacity scaling $L^* = P/(2N) + 1$.
\end{enumerate}
\end{theorem}

\begin{proof}
Identical to Theorem~\ref{thm:linear}, replacing $X_1$ with $\Phi$ and $N$ with $M$. Part (iii) follows from $M = P/(2N)$.
\end{proof}

\begin{remark}[Pointwise nonlinearity does not help in isolation]
If $F_\theta(\mathbf{x}) = \sigma(W\mathbf{x})$ with $W \in \mathbb{R}^{N \times N}$ and $\sigma$ invertible pointwise, then one-shot learning reduces to $W = \sigma^{-1}(X_2) X_1^+$, which still requires $L - 1 \leq N$. Pointwise nonlinearity only helps when applied to an \emph{expanded} representation ($M > N$).
\end{remark}

\begin{theorem}[Fundamental Capacity--Parameter Tradeoff]\label{thm:tradeoff}
For any one-shot sequence memory with $P$ learnable parameters (fixed independently of $L$) and $N$-dimensional patterns:
\[
    L \leq \frac{P}{N} + 1.
\]
\end{theorem}

\begin{proof}[Proof sketch]
Perfect recall of a length-$L$ sequence imposes $(L-1) \times N$ scalar equality constraints (one per component per transition). For a generic parameterization, these constraints are independent, so a solution requires at least $(L-1)N$ free parameters: $P \geq (L-1)N$.
\end{proof}

\begin{corollary}
For a fully connected recurrent network ($P = N^2$): $L \leq N + 1$. This matches the linear pseudoinverse bound exactly, showing the linear method is \textbf{optimal} for this parameter budget.
\end{corollary}

\begin{remark}[Modern Hopfield escapes the bound]
The Modern Hopfield network (Architecture C) uses $P = 2NL$ parameters that \emph{grow with $L$}. Theorem~\ref{thm:tradeoff} does not apply. Its capacity is instead limited by retrieval fidelity: the softmax must correctly identify the current pattern among all stored keys. Ramsauer et al.\ showed this succeeds for up to $L = 2^{cN}$ random patterns on the unit sphere, where $c > 0$ depends on the minimum separation.
\end{remark}

\begin{theorem}[Information Capacity under Energy Constraints]\label{thm:info}
For any sequence memory achieving length $L$ with $N$-dimensional patterns drawn i.i.d.\ from distribution $\mathcal{D}$, the total stored information is $I_{\mathrm{total}} = L \cdot h(\mathcal{D})$, where $h(\mathcal{D})$ is the differential entropy of $\mathcal{D}$.
\begin{enumerate}
    \item[(i)] Under an $L_1$ constraint $\mathbb{E}[\|\mathbf{x}\|_1 / N] \leq c$:
    \[
        h(\mathcal{D}) \leq N\bigl(1 + \ln(2c)\bigr),
    \]
    with equality for i.i.d.\ $\mathrm{Laplace}(0, c)$ components.
    \item[(ii)] Under an $L_2$ constraint $\mathbb{E}[\|\mathbf{x}\|^2 / N] \leq \sigma^2$:
    \[
        h(\mathcal{D}) \leq \frac{N}{2}\ln(2\pi e \sigma^2),
    \]
    with equality for i.i.d.\ $\mathcal{N}(0, \sigma^2)$ components.
\end{enumerate}
\end{theorem}

\begin{proof}
These are standard maximum-entropy results. Under the constraint $\mathbb{E}[|X|] = c$, the Laplace distribution maximizes differential entropy among all distributions on $\mathbb{R}$. Under $\mathbb{E}[X^2] = \sigma^2$, the Gaussian maximizes entropy. The i.i.d.\ structure gives $h(\mathcal{D}) = N \cdot h(D_1)$ where $D_1$ is the marginal.
\end{proof}

\begin{observation}[Separation of capacity and coding]\label{obs:separation}
The sequence length $L$ depends on the architecture and parameter count, \textbf{not} on the pattern distribution (for generic continuous distributions). The information per pattern depends on the distribution, not the architecture. These are independent design choices:
\begin{itemize}
    \item To maximize $L$: choose an architecture with many parameters (e.g., Hopfield).
    \item To maximize information per pattern: choose the entropy-maximizing distribution for the relevant energy constraint (Laplace for $L_1$, Gaussian for $L_2$).
\end{itemize}
\end{observation}

\subsection{Experimental Results}

All experiments used tolerance $\varepsilon = 10^{-4}$ for ``perfect'' recall (per-pattern MSE $< \varepsilon$).

\subsubsection{Experiment 1: Linear Pseudoinverse}

\begin{table}[h]
\centering
\caption{Linear pseudoinverse capacity. $L^*$ = max sequence length, $L^*/N$ = normalized capacity.}
\begin{tabular}{r c c c c}
\toprule
$N$ & Gaussian & Laplace & Sparse (10\%) & Orthogonal \\
\midrule
10  & 11 (1.10) & 11 (1.10) & 5 (0.50) & 11 (1.10) \\
25  & 26 (1.04) & 26 (1.04) & 18 (0.72) & 26 (1.04) \\
50  & 51 (1.02) & 51 (1.02) & 50 (1.00) & 51 (1.02) \\
100 & 100 (1.00) & 100 (1.00) & 99 (0.99) & 101 (1.01) \\
\bottomrule
\end{tabular}
\label{tab:linear}
\end{table}

\textbf{Finding:} Capacity $L^* = N + 1$ confirmed for Gaussian, Laplace, and Orthogonal patterns, exactly matching Theorem~\ref{thm:linear}. Sparse patterns underperform for small $N$ because with 10\% sparsity and $N=10$, each pattern has only 1 nonzero entry, making linear independence fragile.

\subsubsection{Experiment 2: Echo State Network}

\begin{table}[h]
\centering
\caption{Echo state (tanh activation) capacity. $M$ = hidden dimension.}
\begin{tabular}{r r r r r}
\toprule
$N$ & $M = N$ & $M = 2N$ & $M = 4N$ & $M = 8N$ \\
\midrule
25  & 24 (0.96) & 37 (0.74) & 53 (0.53) & 73 (0.36) \\
50  & 41 (0.82) & 62 (0.62) & 86 (0.43) & 123 (0.31) \\
100 & 71 (0.71) & 102 (0.51) & 143 (0.36) & 216 (0.27) \\
\bottomrule
\end{tabular}
\label{tab:echo}
\end{table}

Values shown as $L^*$ ($L^*/M$). \textbf{Finding:} Capacity scales sublinearly with $M$, achieving only $L^*/M \approx 0.3$--$0.9$. The tanh saturation reduces the effective feature dimension: when $|W_1 \mathbf{x}|$ is large, $\tanh$ compresses to $\pm 1$, collapsing distinct features. Gaussian and Laplace performed nearly identically.

\subsubsection{Experiment 3: Modern Hopfield}

We ran two rounds of Hopfield experiments: first with raw (unnormalized) patterns, then with patterns normalized to unit norm.

\begin{table}[h]
\centering
\caption{Modern Hopfield capacity with \textbf{normalized} patterns ($\beta = 50$).}
\begin{tabular}{r r r r}
\toprule
$N$ & Gaussian & Laplace & Orthogonal \\
\midrule
10  & 193 (19.3$N$) & 100 (10.0$N$) & 187 (18.7$N$) \\
25  & $\geq$2000 (80$N$) & $\geq$2000 (80$N$) & $\geq$2000 (80$N$) \\
50  & $\geq$3000 (60$N$) & $\geq$3000 (60$N$) & $\geq$3000 (60$N$) \\
\bottomrule
\end{tabular}
\label{tab:hopfield_norm}
\end{table}

\begin{table}[h]
\centering
\caption{Effect of normalization on capacity ($N = 50$, $\beta = 50$).}
\begin{tabular}{l r r}
\toprule
Configuration & $L^*$ & $L^*/N$ \\
\midrule
Gaussian, normalized & $\geq 3000$ & $\geq 60$ \\
Gaussian, raw & $\geq 3000$ & $\geq 60$ \\
Laplace, normalized & $\geq 3000$ & $\geq 60$ \\
Laplace, raw & 898 & 18.0 \\
\bottomrule
\end{tabular}
\label{tab:normalization}
\end{table}

\textbf{Key findings:}
\begin{enumerate}
    \item The Modern Hopfield network achieves capacity at least $60N$ for $N \geq 50$, dwarfing the linear bound of $N + 1$. True capacity likely much higher (search ceiling hit).
    \item \textbf{Normalization eliminates Laplace's disadvantage.} Raw Laplace patterns achieve only $18N$ capacity, but normalized Laplace matches Gaussian at $\geq 60N$. The problem with raw Laplace is norm variance: patterns with extreme norms dominate the softmax regardless of angular similarity.
    \item For small $N = 10$, Laplace (normalized) still underperforms Gaussian ($10N$ vs $19N$) due to higher mutual coherence of the angular distribution (Laplace concentrates near coordinate axes).
\end{enumerate}

\subsubsection{Experiment 4: Error Accumulation}

For $N = 50$ with the linear pseudoinverse:
\begin{itemize}
    \item $L \leq N$: MSE $\sim 10^{-28}$ (machine precision).
    \item $L = N + 1$: MSE $\sim 10^{-22}$ (still excellent, at theoretical limit).
    \item $L > N + 1$: MSE $\approx 1$ (random-level error). The transition is \textbf{sharp}.
\end{itemize}

\textbf{Notable:} At $L = N + 1$, Laplace patterns showed numerical instability (MSE up to $10^{20}$ in one trial) due to ill-conditioned pattern matrices. This is a practical disadvantage of Laplace coding despite its information-theoretic optimality.

\subsubsection{Experiment 5: Laplace Scale Sweep}

For $N = 50$, varying the Laplace scale parameter $b$: capacity was $L = 51$ for \emph{all} scales tested ($b = 0.1$ to $5.0$). The scale affects only the information per pattern (entropy $= 1 + \ln(2b)$ nats per component), not the sequence length. Total information increases monotonically with scale, from $-1554$ nats ($b = 0.1$, negative entropy = very concentrated) to $8422$ nats ($b = 5.0$).

\section{Connections}

\subsection{Attention is Sequence Memory (Structural)}

The Modern Hopfield retrieval rule $\mathbf{x}_{k+1} = X_2 \operatorname{softmax}(\beta X_1^\top \mathbf{x}_k)$ is \emph{exactly} single-head key-value attention. This is not a metaphor --- it is a mathematical identity established by Ramsauer et al.\ (2020). The sequence memory problem is thus equivalent to: ``How many key-value pairs can a single attention head retrieve reliably?'' The answer is $2^{\Theta(N)}$.

This suggests that transformers' in-context learning ability is, at its core, a sequence memory mechanism. The ``context window'' of a transformer is a Hopfield memory with exponential capacity per attention head.

\subsection{Parametric vs.\ Non-parametric Memory (Structural)}

A fundamental dichotomy emerges:

\begin{center}
\begin{tabular}{lll}
\toprule
& \textbf{Parametric} & \textbf{Non-parametric} \\
\midrule
Example & Linear, Echo State & Modern Hopfield \\
Parameters & Fixed $P$ & $P$ grows with $L$ \\
Capacity & $L \leq P/N + 1$ & $L = 2^{\Theta(N)}$ \\
One-shot rule & Pseudoinverse & Store patterns directly \\
Bottleneck & Parameter count & Retrieval fidelity \\
\bottomrule
\end{tabular}
\end{center}

For $N = 100$ neurons:
\begin{itemize}
    \item Parametric (fully connected, $P = N^2 = 10{,}000$): $L \leq 101$.
    \item Non-parametric (Hopfield): $L$ up to $\sim 2^{50}$, but requires $\sim 2^{50} \times 200$ storage.
\end{itemize}

\subsection{Laplace Optimality: Real but Requires Normalization (Structural)}

The Laplace connection is \emph{information-theoretically} real: under metabolic ($L_1$) constraints, Laplace-distributed activations maximize entropy per pattern. The practical picture is nuanced:

\begin{enumerate}
    \item Sequence capacity $L$ is architecture-determined, independent of distribution.
    \item Raw Laplace patterns degrade Hopfield capacity ($18N$ vs $60N$+) due to norm variance.
    \item \textbf{Normalization eliminates this disadvantage completely}: normalized Laplace matches Gaussian capacity.
    \item After normalization, information content is in the angular distribution; Laplace's entropy advantage applies to the pre-normalization representation.
\end{enumerate}

\textbf{Practical recipe:} Use Laplace-generated patterns (maximizing information under $L_1$ constraints), then normalize to unit norm before Hopfield storage. This gives the best of both worlds: high information density and high retrieval capacity.

The biological parallel: neurons may use Laplace-like marginals for maximal per-neuron information while employing pattern-separation mechanisms (like the hippocampal dentate gyrus) to reduce inter-pattern coherence --- analogous to our normalization step.

\subsection{Orthogonality as Low Coherence (Structural)}

The intuition that ``different neurons should be responsible for different items'' maps precisely to the concept of \emph{mutual coherence}:
\[
    \mu(X) = \max_{i \neq j} \frac{|\mathbf{x}_i^\top \mathbf{x}_j|}{\|\mathbf{x}_i\| \|\mathbf{x}_j\|}.
\]

Measured coherence for $N = 50$, $L = 200$ normalized patterns:

\begin{center}
\begin{tabular}{l c c}
\toprule
Distribution & Max coherence & Avg coherence \\
\midrule
Gaussian & 0.5573 & 0.1133 \\
Laplace  & 0.5796 & 0.1126 \\
Orthogonal & 0.5584 & 0.0855 \\
\bottomrule
\end{tabular}
\end{center}

After normalization, Laplace's coherence disadvantage is small ($\sim$4\% higher max coherence). The average coherence is nearly identical. This explains why normalized Laplace achieves comparable Hopfield capacity to Gaussian for $N \geq 25$.

Lower coherence $\Rightarrow$ better softmax discrimination $\Rightarrow$ higher Hopfield capacity. For the linear pseudoinverse, coherence doesn't affect capacity (any linearly independent set works), but it affects \emph{numerical stability}.

\section{Open Questions}

\begin{question}[Optimal activation distribution for Hopfield]
What distribution on the unit sphere $S^{N-1}$ maximizes Hopfield sequential recall capacity? The uniform distribution (from normalized Gaussian) is a strong candidate, but structured distributions (e.g., designs from coding theory) might do better.
\end{question}

\begin{question}[Breaking the parametric barrier efficiently]
Can a network with $O(N^2)$ parameters and nonlinear dynamics store sequences of length $\omega(N)$? Multi-step dynamics (iterated maps) might achieve this, but one-shot computation of the weights becomes a nonlinear optimization problem.
\end{question}

\begin{question}[Combined Laplace + orthogonality]
Is there a distribution that has Laplace-like marginals (maximizing per-neuron entropy under $L_1$ constraints) while maintaining low mutual coherence (enabling long Hopfield sequences)? This would reconcile the biological observation with the capacity requirements.
\end{question}

\begin{question}[Biological implementation of Hopfield sequence memory]
The hippocampal model in~\cite{plos2024oneshot} achieves $\sim N$ capacity (parametric regime). Does the hippocampus implement something closer to Hopfield's exponential regime, perhaps via pattern separation + associative retrieval?
\end{question}

\section{Assessment}

\textbf{What worked:}
\begin{itemize}
    \item The theoretical framework cleanly separates \emph{sequence capacity} (architecture-dependent) from \emph{information density} (distribution-dependent). Theorems~\ref{thm:linear}--\ref{thm:info} are rigorous and experimentally confirmed.
    \item The parametric/non-parametric dichotomy (Theorem~\ref{thm:tradeoff} vs.\ Hopfield) reveals the fundamental design tradeoff.
    \item Experiments confirm the linear bound $L = N + 1$ is tight.
\end{itemize}

\textbf{What was surprising:}
\begin{itemize}
    \item Raw Laplace patterns are \emph{much worse} for Hopfield ($18N$) --- but \textbf{normalizing them completely eliminates the disadvantage} ($\geq 60N$). The issue was norm variance, not angular distribution.
    \item The echo state tanh network achieves only 27--96\% of its theoretical capacity ($M + 1$), due to activation saturation.
    \item The capacity transition at $L = N + 1$ is extremely sharp: MSE jumps from $10^{-22}$ to $\sim 1$ in a single step.
    \item Hopfield capacity for $N = 50$ exceeds $3000$ ($60N$), far beyond the linear pseudoinverse's $51$ ($\sim N$).
\end{itemize}

\textbf{What was inconclusive:}
\begin{itemize}
    \item The exact Hopfield capacity for $N \geq 50$ remains unknown (exceeded search range of $3000$). Theory predicts $2^{\Theta(N)}$ but the practical crossover to exponential regime needs more exploration.
    \item Whether the parametric barrier $L \leq P/N + 1$ can be broken with nonlinear multi-step dynamics remains open.
\end{itemize}

\textbf{Next steps:}
\begin{itemize}
    \item Push Hopfield capacity search to $10^5$--$10^6$ for $N = 50$--$100$ to find the true ceiling.
    \item Explore coding-theoretic constructions (spherical codes, Grassmannian packings) for the pattern distribution.
    \item Test ReLU echo states (non-saturating) to see if feature expansion efficiency improves.
    \item Investigate whether a hybrid architecture (fixed-size recurrent + external pattern store) can combine the advantages of both regimes.
\end{itemize}

\begin{thebibliography}{9}
\bibitem{plos2024oneshot}
``A neural network model for online one-shot storage of pattern sequences,'' \emph{PLOS One}, 2024.

\bibitem{ramsauer2020hopfield}
H.~Ramsauer et al., ``Hopfield Networks is All You Need,'' \emph{ICLR}, 2021.

\bibitem{neurips2024hopfield}
``Provably Optimal Memory Capacity for Modern Hopfield Models,'' \emph{NeurIPS}, 2024.

\bibitem{matrix_rnn2021}
``Memory Capacity of Recurrent Neural Networks with Matrix Representation,'' \emph{arXiv:2104.07454}, 2021.

\bibitem{sparse_coding}
``Sparse Coding,'' \emph{Scholarpedia}, and related work on $L_1$ sparse coding in V1.
\end{thebibliography}

\end{document}
